\section{Pilot Objectives and Non-Objectives}
\label{sec:objectives}

\subsection{Primary Objectives}

The DUALEXIS pilot programme aims to evaluate---under ethics-approved, governed
conditions---whether a privacy-preserving semantic orchestration architecture can:

\begin{enumerate}
  \item support \textbf{operator decision support} for school safety staff during nominal
        and stress scenarios (crowd movement, exit impairment, multimodal conflict);
  \item integrate with existing \textbf{emergency preparedness protocols} (lockdown,
        evacuation, shelter-in-place, verification) without replacing institutional authority;
  \item maintain \textbf{privacy-preserving situational awareness} via zone-level semantic
        events rather than persistent raw-media archives;
  \item demonstrate \textbf{local processing} on segmented school networks with auditable
        structured-event egress only;
  \item produce evidence for \textbf{safety management system (SMS)} integration discussions
        and future research publications (including Horizon Europe proposals).
\end{enumerate}

\subsection{Secondary Objectives}

\begin{itemize}[noitemsep]
  \item Confirm infrastructure sizing adequacy with school IT for planned zone coverage.
  \item Measure staff comprehension, review workload, and protocol mapping quality in
        tabletop and limited operational exercises.
  \item Document governance artifacts (DPIA inputs, consent/transparency materials, audit logs).
  \item Establish reproducible evaluation baselines for multi-school extension.
\end{itemize}

\subsection{Explicit Non-Objectives}

The pilot does \textbf{not} aim to:
\begin{itemize}[noitemsep]
  \item identify, track, or profile individuals (students, staff, or visitors);
  \item perform facial recognition, voiceprint identification, or biometric analytics;
  \item autonomously trigger lockdowns, dispatch emergency services, or enforce disciplinary action;
  \item replace existing emergency communication systems (PA, radio, municipal services);
  \item claim regulatory certification or universal compliance without institution-specific review;
  \item demonstrate incident-prevention or harm-reduction outcomes without appropriately
        powered studies.
\end{itemize}

\subsection{Success Criteria (Pilot Planning Level)}

Success at the requirements stage means institutional readiness to proceed to governed
pilot execution, defined as:

\begin{itemize}[noitemsep]
  \item signed governance charter and protocol mapping workshop completed;
  \item network and hardware baseline approved by school IT;
  \item DPIA process initiated with documented risk mitigations;
  \item staff training curriculum delivered to safety officers and administrators;
  \item evaluation metrics and data collection plan approved by ethics committee where required.
\end{itemize}

Operational effectiveness thresholds (e.g., situational-awareness scores, false-alert rates)
will be specified in the evaluation plan (Section~\ref{sec:evaluation-metrics}) and
confirmed before inferential analysis.
